% Options for packages loaded elsewhere
\PassOptionsToPackage{unicode}{hyperref}
\PassOptionsToPackage{hyphens}{url}
\documentclass[
]{article}
\usepackage{xcolor}
\usepackage{amsmath,amssymb}
\setcounter{secnumdepth}{-\maxdimen} % remove section numbering
\usepackage{iftex}
\ifPDFTeX
  \usepackage[T1]{fontenc}
  \usepackage[utf8]{inputenc}
  \usepackage{textcomp} % provide euro and other symbols
\else % if luatex or xetex
  \usepackage{unicode-math} % this also loads fontspec
  \defaultfontfeatures{Scale=MatchLowercase}
  \defaultfontfeatures[\rmfamily]{Ligatures=TeX,Scale=1}
\fi
\usepackage{lmodern}
\ifPDFTeX\else
  % xetex/luatex font selection
\fi
% Use upquote if available, for straight quotes in verbatim environments
\IfFileExists{upquote.sty}{\usepackage{upquote}}{}
\IfFileExists{microtype.sty}{% use microtype if available
  \usepackage[]{microtype}
  \UseMicrotypeSet[protrusion]{basicmath} % disable protrusion for tt fonts
}{}
\makeatletter
\@ifundefined{KOMAClassName}{% if non-KOMA class
  \IfFileExists{parskip.sty}{%
    \usepackage{parskip}
  }{% else
    \setlength{\parindent}{0pt}
    \setlength{\parskip}{6pt plus 2pt minus 1pt}}
}{% if KOMA class
  \KOMAoptions{parskip=half}}
\makeatother
\usepackage{color}
\usepackage{fancyvrb}
\newcommand{\VerbBar}{|}
\newcommand{\VERB}{\Verb[commandchars=\\\{\}]}
\DefineVerbatimEnvironment{Highlighting}{Verbatim}{commandchars=\\\{\}}
% Add ',fontsize=\small' for more characters per line
\newenvironment{Shaded}{}{}
\newcommand{\AlertTok}[1]{\textcolor[rgb]{1.00,0.00,0.00}{\textbf{#1}}}
\newcommand{\AnnotationTok}[1]{\textcolor[rgb]{0.38,0.63,0.69}{\textbf{\textit{#1}}}}
\newcommand{\AttributeTok}[1]{\textcolor[rgb]{0.49,0.56,0.16}{#1}}
\newcommand{\BaseNTok}[1]{\textcolor[rgb]{0.25,0.63,0.44}{#1}}
\newcommand{\BuiltInTok}[1]{\textcolor[rgb]{0.00,0.50,0.00}{#1}}
\newcommand{\CharTok}[1]{\textcolor[rgb]{0.25,0.44,0.63}{#1}}
\newcommand{\CommentTok}[1]{\textcolor[rgb]{0.38,0.63,0.69}{\textit{#1}}}
\newcommand{\CommentVarTok}[1]{\textcolor[rgb]{0.38,0.63,0.69}{\textbf{\textit{#1}}}}
\newcommand{\ConstantTok}[1]{\textcolor[rgb]{0.53,0.00,0.00}{#1}}
\newcommand{\ControlFlowTok}[1]{\textcolor[rgb]{0.00,0.44,0.13}{\textbf{#1}}}
\newcommand{\DataTypeTok}[1]{\textcolor[rgb]{0.56,0.13,0.00}{#1}}
\newcommand{\DecValTok}[1]{\textcolor[rgb]{0.25,0.63,0.44}{#1}}
\newcommand{\DocumentationTok}[1]{\textcolor[rgb]{0.73,0.13,0.13}{\textit{#1}}}
\newcommand{\ErrorTok}[1]{\textcolor[rgb]{1.00,0.00,0.00}{\textbf{#1}}}
\newcommand{\ExtensionTok}[1]{#1}
\newcommand{\FloatTok}[1]{\textcolor[rgb]{0.25,0.63,0.44}{#1}}
\newcommand{\FunctionTok}[1]{\textcolor[rgb]{0.02,0.16,0.49}{#1}}
\newcommand{\ImportTok}[1]{\textcolor[rgb]{0.00,0.50,0.00}{\textbf{#1}}}
\newcommand{\InformationTok}[1]{\textcolor[rgb]{0.38,0.63,0.69}{\textbf{\textit{#1}}}}
\newcommand{\KeywordTok}[1]{\textcolor[rgb]{0.00,0.44,0.13}{\textbf{#1}}}
\newcommand{\NormalTok}[1]{#1}
\newcommand{\OperatorTok}[1]{\textcolor[rgb]{0.40,0.40,0.40}{#1}}
\newcommand{\OtherTok}[1]{\textcolor[rgb]{0.00,0.44,0.13}{#1}}
\newcommand{\PreprocessorTok}[1]{\textcolor[rgb]{0.74,0.48,0.00}{#1}}
\newcommand{\RegionMarkerTok}[1]{#1}
\newcommand{\SpecialCharTok}[1]{\textcolor[rgb]{0.25,0.44,0.63}{#1}}
\newcommand{\SpecialStringTok}[1]{\textcolor[rgb]{0.73,0.40,0.53}{#1}}
\newcommand{\StringTok}[1]{\textcolor[rgb]{0.25,0.44,0.63}{#1}}
\newcommand{\VariableTok}[1]{\textcolor[rgb]{0.10,0.09,0.49}{#1}}
\newcommand{\VerbatimStringTok}[1]{\textcolor[rgb]{0.25,0.44,0.63}{#1}}
\newcommand{\WarningTok}[1]{\textcolor[rgb]{0.38,0.63,0.69}{\textbf{\textit{#1}}}}
\usepackage{longtable,booktabs,array}
\usepackage{calc} % for calculating minipage widths
% Correct order of tables after \paragraph or \subparagraph
\usepackage{etoolbox}
\makeatletter
\patchcmd\longtable{\par}{\if@noskipsec\mbox{}\fi\par}{}{}
\makeatother
% Allow footnotes in longtable head/foot
\IfFileExists{footnotehyper.sty}{\usepackage{footnotehyper}}{\usepackage{footnote}}
\makesavenoteenv{longtable}
\setlength{\emergencystretch}{3em} % prevent overfull lines
\providecommand{\tightlist}{%
  \setlength{\itemsep}{0pt}\setlength{\parskip}{0pt}}
\usepackage{bookmark}
\IfFileExists{xurl.sty}{\usepackage{xurl}}{} % add URL line breaks if available
\urlstyle{same}
\hypersetup{
  hidelinks,
  pdfcreator={LaTeX via pandoc}}

\author{}
\date{}

\begin{document}

\section{Project: The Classical Resonance Engine (2-Bit / 4-State
Demonstrator)}\label{project-the-classical-resonance-engine-2-bit-4-state-demonstrator}

\subsubsection{Demystifying ``Quantum'' Superposition and Oracle Search
Using Room-Temperature Analog
Waves}\label{demystifying-quantum-superposition-and-oracle-search-using-room-temperature-analog-waves}

This open-source hardware project implements Grover's search algorithm
using high-frequency, room-temperature operational amplifiers. By
bypassing abstract matrix formalism and cryogenic dependencies, this
circuit models the search space as a parallel network of phase-locked
oscillators. It demonstrates that ``quantum'' amplitude amplification is
structurally isomorphic to classical wave mechanics,
constructive/destructive interference, and analog signal mixing.

\begin{center}\rule{0.5\linewidth}{0.5pt}\end{center}

\subsection{1. Algorithmic Architecture \& Engineering
Theory}\label{algorithmic-architecture-engineering-theory}

\subsubsection{The Classical Phase View}\label{the-classical-phase-view}

In academic quantum computing literature, Grover's search maps an
unstructured database to an abstract, unobservable Hilbert space where
states exist in a fragile superposition. The state vector is rotated
blindly for \(R \approx \frac{\pi}{4}\sqrt{N}\) steps until a
destructive measurement collapses the vector.

This circuit strips away that vocabulary. We represent a 2-bit database
(\(N = 2^2 = 4\) states) as four physical, concurrent AC voltage signals
oscillating at a fixed carrier frequency
(\(f_0 \approx 16\text{ kHz}\)).

\begin{itemize}
\tightlist
\item
  \textbf{Superposition = Uniform Wave Distribution:} A single reference
  clock pumps identical, in-phase signals across four independent
  operational amplifier buffers. Every state is active simultaneously.
\item
  \textbf{The Oracle = \(180^\circ\) (\(\pi\) Radian) Phase Shift:}
  Selecting a ``target'' coordinate via a mechanical switch routes that
  specific channel through a unity-gain inverting amplifier
  (\(V_{\text{out}} = -V_{\text{in}}\)). It flips the sign of that
  coordinate's amplitude.
\item
  \textbf{The Diffusion Mixer = DC Offset Inversion:} A central
  inverting summing amplifier continuously adds the signals from all
  four nodes. This calculates the global arithmetic mean (\(\mu\)). By
  feeding the inverted mean back into the array, the circuit performs an
  instantaneous, continuous \textbf{Inversion About the Average}:
\end{itemize}

\[V_{\text{out\_channel}} = 2V_{\text{mean}} - V_{\text{in\_channel}}\]

\subsubsection{Continuous Early Termination (Why Analog Beats the
Textbooks)}\label{continuous-early-termination-why-analog-beats-the-textbooks}

Because these are macroscopic, classical voltage signals, there is no
``wavefunction collapse.'' A serious hobbyist can use a standard
oscilloscope probe or a window comparator to monitor the nodes
continuously without altering the data trajectory.

When a target's phase is flipped, the continuous active feedback loop
causes the target channel's voltage to spike past the baseline almost
instantly, while the incorrect channels suffer destructive interference
and dim out. A comparator detects this threshold breakout and halts
immediately---physically demonstrating \textbf{numerical early
termination} at room temperature.

\begin{center}\rule{0.5\linewidth}{0.5pt}\end{center}

\subsection{2. Complete Bill of Materials
(BOM)}\label{complete-bill-of-materials-bom}

All components listed are inexpensive, generic, and readily available
from major distributors (Mouser, DigiKey, JameCo) or standard
electronics hobby bins.

\subsubsection{Active Semiconductors}\label{active-semiconductors}

\begin{longtable}[]{@{}
  >{\raggedright\arraybackslash}p{(\linewidth - 8\tabcolsep) * \real{0.2000}}
  >{\raggedright\arraybackslash}p{(\linewidth - 8\tabcolsep) * \real{0.2000}}
  >{\raggedright\arraybackslash}p{(\linewidth - 8\tabcolsep) * \real{0.2000}}
  >{\raggedright\arraybackslash}p{(\linewidth - 8\tabcolsep) * \real{0.2000}}
  >{\raggedright\arraybackslash}p{(\linewidth - 8\tabcolsep) * \real{0.2000}}@{}}
\toprule\noalign{}
\begin{minipage}[b]{\linewidth}\raggedright
Qty
\end{minipage} & \begin{minipage}[b]{\linewidth}\raggedright
Component Type
\end{minipage} & \begin{minipage}[b]{\linewidth}\raggedright
Part Number
\end{minipage} & \begin{minipage}[b]{\linewidth}\raggedright
Description
\end{minipage} & \begin{minipage}[b]{\linewidth}\raggedright
Purpose
\end{minipage} \\
\midrule\noalign{}
\endhead
\bottomrule\noalign{}
\endlastfoot
2 & Quad Operational Amplifier & \textbf{TL074CN} (or TL084) & High
Slew-Rate JFET-Input Quad Op-Amp & Array buffers, phase inverters, and
summing mixers. \\
1 & Quad Comparator & \textbf{LM339N} & Low-Voltage Open-Collector Quad
Comparator & Direct window-threshold discriminator for early
breakout. \\
\end{longtable}

\subsubsection{Passives \& Electromechanical (1/4W, 1\% Tolerance
Recommended)}\label{passives-electromechanical-14w-1-tolerance-recommended}

\begin{longtable}[]{@{}
  >{\raggedright\arraybackslash}p{(\linewidth - 6\tabcolsep) * \real{0.2500}}
  >{\raggedright\arraybackslash}p{(\linewidth - 6\tabcolsep) * \real{0.2500}}
  >{\raggedright\arraybackslash}p{(\linewidth - 6\tabcolsep) * \real{0.2500}}
  >{\raggedright\arraybackslash}p{(\linewidth - 6\tabcolsep) * \real{0.2500}}@{}}
\toprule\noalign{}
\begin{minipage}[b]{\linewidth}\raggedright
Qty
\end{minipage} & \begin{minipage}[b]{\linewidth}\raggedright
Value / Part
\end{minipage} & \begin{minipage}[b]{\linewidth}\raggedright
Type
\end{minipage} & \begin{minipage}[b]{\linewidth}\raggedright
Purpose / Placement
\end{minipage} \\
\midrule\noalign{}
\endhead
\bottomrule\noalign{}
\endlastfoot
12 & \textbf{10 kΩ} & Metal Film Resistor & Crucial for matched
unity-gain inversion and summing nodes. \\
4 & \textbf{1 kΩ} & Carbon Film Resistor & Current-limiting resistors
for the output display LEDs. \\
4 & \textbf{1N4148} & Fast Switching Diode & Half-wave rectifiers to
convert AC wave amplitude to DC levels. \\
1 & \textbf{10 nF} (\(0.01\mu\text{F}\)) & Ceramic Disc Capacitor & DC
blocking / low-pass filter smoothing for envelope detectors. \\
1 & \textbf{4-Position DIP} & SPST Switch Block & Represents the 2-bit
search database inputs (States 00, 01, 10, 11). \\
4 & \textbf{Standard LED} & Red or Green T-1 3/4 & Visual amplitude
resonance indicator per coordinate lane. \\
\end{longtable}

\begin{center}\rule{0.5\linewidth}{0.5pt}\end{center}

\subsection{3. LTspice Reference Circuit
Architecture}\label{ltspice-reference-circuit-architecture}

To verify the circuit's phase-mixing operations before modifying real
hardware, copy the raw netlist text block below and save it locally as
\texttt{grover\_resonance.net}. You can open this file directly inside
LTspice to plot the continuous wave amplification.

\begin{Shaded}
\begin{Highlighting}[]
\NormalTok{* Classical Grover Resonance Engine {-} 4{-}State Analog Netlist Snapshot}
\NormalTok{.OPTIONS plotwinsize=0 numdgt=7}
\NormalTok{.PARAM FREQ=15915 RMATCH=10k}

\NormalTok{* 1. Carrier Signal Reference Generator (15.9 kHz)}
\NormalTok{V\_Ref Ref\_CLK 0 SINE(0 1 \{FREQ\})}
\NormalTok{R\_Ref\_Load Ref\_CLK 0 100k}

\NormalTok{* 2. State Channel Array Buffers (Non{-}Inverting Configuration)}
\NormalTok{X\_Buf0 Ref\_CLK Node\_0\_In VCC VEE TL074}
\NormalTok{X\_Buf1 Ref\_CLK Node\_1\_In VCC VEE TL074}
\NormalTok{X\_Buf2 Ref\_CLK Node\_2\_In VCC VEE TL074}
\NormalTok{X\_Buf3 Ref\_CLK Node\_3\_In VCC VEE TL074}

\NormalTok{* 3. Mechanical Oracle Matrix Simulation (State 2 / Node 2 Phase Flipped)}
\NormalTok{* To switch targets in LTspice, manually change which node connects to the inverter.}
\NormalTok{X\_Inv2 Node\_2\_In Node\_2\_Oracle VCC VEE TL074 Res\_Inv1 Node\_2\_In Node\_2\_Oracle \{RMATCH\}}

\NormalTok{* 4. Central Diffusion Summing Mixer (Computes Mean and Multiplies by 2)}
\NormalTok{R\_Sum0 Node\_0\_In Sum\_Node \{RMATCH\}}
\NormalTok{R\_Sum1 Node\_1\_In Sum\_Node \{RMATCH\}}
\NormalTok{R\_Sum2 Node\_2\_Oracle Sum\_Node \{RMATCH\}}
\NormalTok{R\_Sum3 Node\_3\_In Sum\_Node \{RMATCH\}}

\NormalTok{X\_Mixer Sum\_Node Mean\_Out VCC VEE TL074}
\NormalTok{R\_Feedback Sum\_Node Mean\_Out \{RMATCH\}}

\NormalTok{* 5. Envelope Demodulator \& Early Termination Peak Check}
\NormalTok{D\_Det0 Node\_2\_Oracle Envelope\_0 1N4148}
\NormalTok{C\_Smooth0 Envelope\_0 0 10nF}
\NormalTok{R\_Bleed0 Envelope\_0 0 100k}

\NormalTok{* 6. Power Supply Rails}
\NormalTok{V\_Pos VCC 0 12}
\NormalTok{V\_Neg VEE 0 {-}12}

\NormalTok{.TRAN 0 2m 0 1u}
\NormalTok{.END}
\end{Highlighting}
\end{Shaded}

\begin{center}\rule{0.5\linewidth}{0.5pt}\end{center}

\subsection{4. Hardware Calibration \& Scope
Diagnostics}\label{hardware-calibration-scope-diagnostics}

Once you assemble the circuit on a solderless breadboard, follow this
validation protocol to document your performance metrics:

\subsubsection{Step A: Power Configuration \& Baseline
Balance}\label{step-a-power-configuration-baseline-balance}

\begin{enumerate}
\def\labelenumi{\arabic{enumi}.}
\tightlist
\item
  Connect a regulated dual-rail power supply (\(\pm12\text{V}\)DC or
  \(\pm15\text{V}\)DC) to the \texttt{VCC} and \texttt{VEE} pins of the
  TL074 packages. Decouple each rail using a \(0.1\mu\text{F}\) ceramic
  capacitor placed immediately adjacent to the chip pins to kill
  high-frequency switching noise.
\item
  Turn \textbf{all DIP switches OFF}. All four monitor LEDs should glow
  with identical, faint, baseline intensity.
\item
  Hook your oscilloscope Probe Channel 1 to \texttt{Mean\_Out}. With no
  target selected, the sum of four perfectly symmetric in-phase waves
  divided by their count yields a clean, unshifted reference sine wave.
\end{enumerate}

\subsubsection{Step B: Oracle Execution \& Signal
Tracking}\label{step-b-oracle-execution-signal-tracking}

\begin{enumerate}
\def\labelenumi{\arabic{enumi}.}
\tightlist
\item
  Flip \textbf{DIP Switch 2 ON} (this marks coordinate state \(10_2\) as
  the target proposition).
\item
  Move your oscilloscope probe to monitor the output pin of the target
  channel. You will see the sine wave immediately shift position by
  exactly half a wavelength (\(180^\circ\) phase inversion).
\item
  Monitor \texttt{Mean\_Out} now. The inversion of a single lane causes
  the global average to drop. The central mixer amplifies this
  discrepancy, pushing a massive feedback current back into the array.
\end{enumerate}

\subsubsection{Step C: Early Termination Threshold
Verification}\label{step-c-early-termination-threshold-verification}

\begin{enumerate}
\def\labelenumi{\arabic{enumi}.}
\tightlist
\item
  Probe the output of your half-wave envelope diode detector
  (\texttt{Envelope\_0}).
\item
  Watch the DC voltage level on your screen. It does not wait for a
  clock cycle; it surges continuously from a baseline of roughly
  \(0.2\text{V}\) up to a sharp peak of over \(1.5\text{V}\) within
  \textbf{two wave periods} (\(\approx 120\text{ microseconds}\)).
\item
  The LM339 comparator catches this rising edge, instantly tripping its
  output pin to open-collector ground. This turns on your ``Target
  Locked'' indicator LED, proving you have caught the correct index
  using pure, room-temperature analog wave mechanics.
\end{enumerate}

\begin{center}\rule{0.5\linewidth}{0.5pt}\end{center}

\subsection{5. LM339 Window Comparator Hardware
Topology}\label{lm339-window-comparator-hardware-topology}

To achieve reliable early termination without processor overhead, we
route the demodulated analog envelope through a hardware \textbf{Window
Comparator}. This ensures the system instantly locks and flags an
interrupt only when a coordinate lane's wave amplitude breaches the
background noise threshold.

\subsubsection{Threshold Reference
Dividers}\label{threshold-reference-dividers}

Connect three \(10\text{ k}\Omega\) resistors in series between the
\(+12\text{V}\) (VCC) rail and System Ground (GND) to establish a
dual-threshold reference network: * \textbf{High Threshold Reference
Point (\(\approx 3.0\text{V}\)):} Taken from the tap between the first
and second resistors. * \textbf{Low Threshold Reference Point
(\(\approx 1.5\text{V}\)):} Taken from the tap between the second and
third resistors.

\subsubsection{Comparator Integration
Map}\label{comparator-integration-map}

\begin{itemize}
\tightlist
\item
  \textbf{LM339 Stage 1 (Over-Voltage Monitor):}

  \begin{itemize}
  \tightlist
  \item
    Connect the \textbf{Inverting Input (-)} to the \textbf{High
    Threshold Reference Point (\(\approx 3.0\text{V}\))}.
  \item
    Connect the \textbf{Non-Inverting Input (+)} to your
    \textbf{Envelope Detector Output}.
  \end{itemize}
\item
  \textbf{LM339 Stage 2 (Under-Voltage Monitor):}

  \begin{itemize}
  \tightlist
  \item
    Connect the \textbf{Inverting Input (-)} to your \textbf{Envelope
    Detector Output}.
  \item
    Connect the \textbf{Non-Inverting Input (+)} to the \textbf{Low
    Threshold Reference Point (\(\approx 1.5\text{V}\))}.
  \end{itemize}
\end{itemize}

\subsubsection{Combined Logic Output}\label{combined-logic-output}

\begin{itemize}
\tightlist
\item
  Tie the Open-Collector Output Pin of Stage 1 directly to the
  Open-Collector Output Pin of Stage 2. This hardwires an analog
  \textbf{Window Comparator Logic Block}.
\item
  Connect a single \(4.7\text{k}\Omega\) pull-up resistor from this
  combined output node up to the Arduino Nano's \(+5\text{V}\) rail.
\item
  Route this same combined output node directly into \textbf{Arduino Pin
  D2} to serve as the hardware interrupt trigger.
\end{itemize}

\begin{center}\rule{0.5\linewidth}{0.5pt}\end{center}

\subsection{6. Arduino Nano Automation Script (The Interface
Layer)}\label{arduino-nano-automation-script-the-interface-layer}

Instead of using physical DIP switches to select database targets, this
script routes the selection lines through digital logic gates (or analog
multiplexer channels like the 4051/74HC4067). It configures a high-speed
hardware interrupt pin (\texttt{D2}) to monitor the LM339 output,
logging the exact execution time down to the microsecond.

\begin{Shaded}
\begin{Highlighting}[]
\CommentTok{/*}
\CommentTok{ *  The Classical Resonance Engine {-} Hybrid Controller Interface}
\CommentTok{ *  Automating Analog Wave Formulations on Standard Silicon Microcontrollers}
\CommentTok{ */}

\PreprocessorTok{\#include }\ImportTok{\textless{}Arduino.h\textgreater{}}

\CommentTok{// Pin Allocation Metrics}
\AttributeTok{const} \DataTypeTok{int}\NormalTok{ TARGET\_SELECT\_A }\OperatorTok{=} \DecValTok{4}\OperatorTok{;} \CommentTok{// Digital bit 0 representing address space}
\AttributeTok{const} \DataTypeTok{int}\NormalTok{ TARGET\_SELECT\_B }\OperatorTok{=} \DecValTok{5}\OperatorTok{;} \CommentTok{// Digital bit 1 representing address space}
\AttributeTok{const} \DataTypeTok{int}\NormalTok{ INTERRUPT\_PIN   }\OperatorTok{=} \DecValTok{2}\OperatorTok{;} \CommentTok{// Dedicated hardware interrupt pin (INT0)}

\CommentTok{// Volatile trackers to communicate safely with the interrupt service routine}
\AttributeTok{volatile} \DataTypeTok{bool}\NormalTok{     resonance\_locked }\OperatorTok{=} \KeywordTok{false}\OperatorTok{;}
\AttributeTok{volatile} \DataTypeTok{uint32\_t}\NormalTok{ stop\_time\_us     }\OperatorTok{=} \DecValTok{0}\OperatorTok{;}
\DataTypeTok{uint32\_t}\NormalTok{          start\_time\_us    }\OperatorTok{=} \DecValTok{0}\OperatorTok{;}

\CommentTok{// Hardware Interrupt Service Routine (ISR)}
\CommentTok{// Triggered instantly the microsecond the LM339 Window Comparator trips}
\DataTypeTok{void}\NormalTok{ ICACHE\_RAM\_ATTR on\_resonance\_lock}\OperatorTok{()} \OperatorTok{\{}
\NormalTok{    stop\_time\_us     }\OperatorTok{=}\NormalTok{ micros}\OperatorTok{();}
\NormalTok{    resonance\_locked }\OperatorTok{=} \KeywordTok{true}\OperatorTok{;}
\OperatorTok{\}}

\DataTypeTok{void}\NormalTok{ setup}\OperatorTok{()} \OperatorTok{\{}
\NormalTok{    Serial}\OperatorTok{.}\NormalTok{begin}\OperatorTok{(}\DecValTok{115200}\OperatorTok{);}
    \ControlFlowTok{while} \OperatorTok{(!}\NormalTok{Serial}\OperatorTok{)} \OperatorTok{\{} \OperatorTok{;} \OperatorTok{\}} \CommentTok{// Wait for terminal initialization}

    \CommentTok{// Configure selection lines to manipulate the analog oracle switches}
\NormalTok{    pinMode}\OperatorTok{(}\NormalTok{TARGET\_SELECT\_A}\OperatorTok{,}\NormalTok{ OUTPUT}\OperatorTok{);}
\NormalTok{    pinMode}\OperatorTok{(}\NormalTok{TARGET\_SELECT\_B}\OperatorTok{,}\NormalTok{ OUTPUT}\OperatorTok{);}
    
    \CommentTok{// Configure interrupt line with internal pull{-}up safety}
\NormalTok{    pinMode}\OperatorTok{(}\NormalTok{INTERRUPT\_PIN}\OperatorTok{,}\NormalTok{ INPUT\_PULLUP}\OperatorTok{);}
    
    \CommentTok{// Attach interrupt handler to catch the falling edge (LM339 open{-}collector drop)}
\NormalTok{    attachInterrupt}\OperatorTok{(}\NormalTok{digitalPinToInterrupt}\OperatorTok{(}\NormalTok{INTERRUPT\_PIN}\OperatorTok{),}\NormalTok{ on\_resonance\_lock}\OperatorTok{,}\NormalTok{ FALLING}\OperatorTok{);}

\NormalTok{    Serial}\OperatorTok{.}\NormalTok{println}\OperatorTok{(}\StringTok{"================================================="}\OperatorTok{);}
\NormalTok{    Serial}\OperatorTok{.}\NormalTok{println}\OperatorTok{(}\StringTok{"   Analog Grover Resonance Engine Initialized   "}\OperatorTok{);}
\NormalTok{    Serial}\OperatorTok{.}\NormalTok{println}\OperatorTok{(}\StringTok{"================================================="}\OperatorTok{);}
\OperatorTok{\}}

\DataTypeTok{void}\NormalTok{ loop}\OperatorTok{()} \OperatorTok{\{}
    \CommentTok{// 1. Reset state flags}
\NormalTok{    resonance\_locked }\OperatorTok{=} \KeywordTok{false}\OperatorTok{;}
    
    \CommentTok{// 2. Select a target index inside the 2{-}bit space (Example: State 10\_2)}
    \CommentTok{// Bit 0 = Low, Bit 1 = High}
\NormalTok{    digitalWrite}\OperatorTok{(}\NormalTok{TARGET\_SELECT\_A}\OperatorTok{,}\NormalTok{ LOW}\OperatorTok{);}
\NormalTok{    digitalWrite}\OperatorTok{(}\NormalTok{TARGET\_SELECT\_B}\OperatorTok{,}\NormalTok{ HIGH}\OperatorTok{);}
    
\NormalTok{    Serial}\OperatorTok{.}\NormalTok{println}\OperatorTok{(}\StringTok{"[System] Phase{-}inversion oracle activated for State 10..."}\OperatorTok{);}

    \CommentTok{// 3. Mark the precise starting epoch}
\NormalTok{    start\_time\_us }\OperatorTok{=}\NormalTok{ micros}\OperatorTok{();}

    \CommentTok{// 4. Force the processor to wait while the analog op{-}amps mix the wave states}
    \CommentTok{// The microcontroller does absolutely nothing here; the analog wires do the calculations.}
    \ControlFlowTok{while} \OperatorTok{(!}\NormalTok{resonance\_locked}\OperatorTok{)} \OperatorTok{\{}
        \CommentTok{// Enforce a simple safety timeout barrier to prevent infinite stall}
        \ControlFlowTok{if} \OperatorTok{(}\NormalTok{micros}\OperatorTok{()} \OperatorTok{{-}}\NormalTok{ start\_time\_us }\OperatorTok{\textgreater{}} \DecValTok{50000}\OperatorTok{)} \OperatorTok{\{} 
\NormalTok{            Serial}\OperatorTok{.}\NormalTok{println}\OperatorTok{(}\StringTok{"[Error] Resonance threshold timeout. Check circuit tuning."}\OperatorTok{);}
            \ControlFlowTok{break}\OperatorTok{;}
        \OperatorTok{\}}
    \OperatorTok{\}}

    \CommentTok{// 5. Output precise benchmark telemetry metrics}
    \ControlFlowTok{if} \OperatorTok{(}\NormalTok{resonance\_locked}\OperatorTok{)} \OperatorTok{\{}
        \DataTypeTok{uint32\_t}\NormalTok{ total\_execution\_time }\OperatorTok{=}\NormalTok{ stop\_time\_us }\OperatorTok{{-}}\NormalTok{ start\_time\_us}\OperatorTok{;}
        
\NormalTok{        Serial}\OperatorTok{.}\NormalTok{println}\OperatorTok{(}\StringTok{"{-}{-}{-}{-}{-}{-}{-}{-}{-}{-}{-}{-}{-}{-}{-}{-}{-}{-}{-}{-}{-}{-}{-}{-}{-}{-}{-}{-}{-}{-}{-}{-}{-}{-}{-}{-}{-}{-}{-}{-}{-}{-}{-}{-}{-}{-}{-}{-}{-}"}\OperatorTok{);}
\NormalTok{        Serial}\OperatorTok{.}\NormalTok{println}\OperatorTok{(}\StringTok{" {-}\textgreater{} SUCCESS: Target Coordinate Locked By Analog Waves!"}\OperatorTok{);}
\NormalTok{        Serial}\OperatorTok{.}\NormalTok{print}\OperatorTok{(}\StringTok{" {-}\textgreater{} Real{-}World Execution Time: "}\OperatorTok{);}
\NormalTok{        Serial}\OperatorTok{.}\NormalTok{print}\OperatorTok{(}\NormalTok{total\_execution\_time}\OperatorTok{);}
\NormalTok{        Serial}\OperatorTok{.}\NormalTok{println}\OperatorTok{(}\StringTok{" microseconds."}\OperatorTok{);}
\NormalTok{        Serial}\OperatorTok{.}\NormalTok{println}\OperatorTok{(}\StringTok{"{-}{-}{-}{-}{-}{-}{-}{-}{-}{-}{-}{-}{-}{-}{-}{-}{-}{-}{-}{-}{-}{-}{-}{-}{-}{-}{-}{-}{-}{-}{-}{-}{-}{-}{-}{-}{-}{-}{-}{-}{-}{-}{-}{-}{-}{-}{-}{-}{-}"}\OperatorTok{);}
    \OperatorTok{\}}

    \CommentTok{// Hold execution for 5 seconds before launching the next search block sweep}
\NormalTok{    delay}\OperatorTok{(}\DecValTok{5000}\OperatorTok{);}
\OperatorTok{\}}
\end{Highlighting}
\end{Shaded}


\end{document}
